% !PRINTED 533
\weberpage{533}

%<<CONT
and consequently $\psi'$ runs through a system of representatives of the discriminant $-4m$ at the same time as $\psi$. But $\omega_1$ runs through the same sequence of values as $\omega_2$, though in a different order, and accordingly summation of formula (30) gives:
\weq{31}{\Pi\,\frac{f(\omega)}{\sqrt[4]2} = 1,}
if the product $\Pi$ refers to the roots with positive imaginary part of a complete system of forms with the discriminant $-4m$.

\begin{quote}
\gesperrt{In what follows we shall, for the sake of brevity, call this sequence of values of the $\omega$ a complete system of roots of the discriminant $-4m$.}
\end{quote}

\S\,142. \gesperrt{The Norms of the Class Invariants $f(\omega)$.}

We let $\omega$ run through a complete system of roots of the discriminant $-4m$, and we assume that in the form $(A, 2B, C)$, whose root is $\omega$, $A, C$ are odd and $A$ is relatively prime to $m$, which involves no further restriction; then by \S\,126 the 24th powers of $f(\omega)$ are class invariants and their norm is a power of 2.

In a two-sided class there is always a representative of one of the two forms:
\[
(A,\,0,\,C),\quad (A,\,2B,\,A),
\]
in which $A$ may be assumed to be odd. In the first case $\omega$ is purely imaginary and consequently [\S\,24, (11)]:
\[
f(\omega),\quad f_1(\omega),\quad f_2(\omega)
\]
are real and positive; in the second case $\omega$ and $1:\omega$ are conjugate imaginary, and consequently:
\[
f(\omega) = f\!\left(-\frac{1}{\omega}\right)
\]
is real and
\[
f_1(\omega) = f_2\!\left(-\frac{1}{\omega}\right),\quad f_2(\omega) = f_1\!\left(-\frac{1}{\omega}\right)
\]
are conjugate imaginary, and by the formula $f(\omega)f_1(\omega)f_2(\omega)=\sqrt 2$ here too $f(\omega)$ is positive.

We may therefore always assume the representative $(A, 2B, C)$ to be so chosen that $A$ and $C$ are odd and that $f(\omega)$ becomes real and positive for a two-sided class; if
%CONT>>
